\section{Algoritmo de Deutsch}\label{sec:deutsch-algorithm}
El algoritmo de Deutsch es uno de los primeros algoritmos cuánticos que resultó más eficiente que su contraparte clásica.

Sea $f:\B\longrightarrow\B$ una función.
Solo existen cuatro funciones distintas: las funciones constantes $0$ y $1$, y las funciones equilibradas identidad y la función opuesta o de inversión.
Si necesitamos saber si $f$ es constante o equilibrada, debemos evaluar $f$ dos veces, pero con un algoritmo cuántico basta con evaluarla una sola vez, como demostró Deutsch~\cite{Deutsch1985}.

La observación clave es que $f$ es constante si y solo si $f(0)\oplus f(1)=0$.
El algoritmo cuántico extrae exactamente este único bit de información con una sola evaluación de $f$, consultando el oráculo sobre una superposición de ambas entradas.

Sea $\gate{U}$ el operador unitario definido sobre la base computacional por
\[
  \gate{U}\ket{xy}=\ket{x(y\oplus f(x))},
\]
donde $\oplus$ es la suma módulo 2.
Si aplicamos el operador $\gate{U}$ sobre el estado $\ket{+-}$ obtenemos:
\[
  \begin{split}
    \gate{U}\ket{+-} & = \frac{1}{\sqrt{2}}\gate{U}(\ket{+0}-\ket{+1})=\frac{1}{2}\gate{U}(\ket{00}+\ket{10}-\ket{01}-\ket{11})= \\
                     & = \frac{1}{2}(\ket{0(0\oplus f(0))}+\ket{1(0\oplus f(1))}-\ket{0(1\oplus f(0))}-\ket{1(1\oplus f(1))})=   \\
                     & =\begin{cases}
                          \frac{1}{2}(\ket{00}+\ket{10}-\ket{01}-\ket{11})=\ket{+-}\text{ si $f(0)=0$ y $f(1) = 0$}  \\
                          \frac{1}{2}(\ket{00}+\ket{11}-\ket{01}-\ket{10})=\ket{--}\text{ si $f(0)=0$ y $f(1) = 1$}  \\
                          \frac{1}{2}(\ket{01}+\ket{10}-\ket{00}-\ket{11})=-\ket{--}\text{ si $f(0)=1$ y $f(1) = 0$} \\
                          \frac{1}{2}(\ket{01}+\ket{11}-\ket{00}-\ket{10})=-\ket{+-}\text{ si $f(0)=1$ y $f(1) = 1$}
                        \end{cases}
  \end{split}
\]
por lo que el estado resultante es
\[
  \gate{U}\ket{+-}=(-1)^{f(0)}\begin{cases}
    \ket{+-} & \text{ si $f$ es constante}   \\
    \ket{--} & \text{ si $f$ es equilibrada}
  \end{cases}.
\]

Consideremos un circuito cuántico que implementa el algoritmo de Deutsch.
El circuito aplica, por este orden, $\gate{X}_2$, $\gate{H}_{1,2}$, $\gate{U}$ y $\gate{H}_1$; con el convenio de Dirac se escribe en orden inverso:
\[
  \gate{D}=\gate{H}_1\gate{U}\gate{H}_{1,2}\gate{X}_2
\]

Aplicando el circuito al estado inicial $\ket{00}$ y usando el cálculo anterior:
\begin{equation}
  \label{eq:deutchs-circuit}
  \begin{split}
    \gate{D}\ket{00} & =\gate{H}_1\gate{U}\gate{H}_{1,2}\gate{X}_2\ket{00}=\gate{H}_1\gate{U}\gate{H}_{1,2}\ket{01}=\gate{H}_1\gate{U}\ket{+-}= \\
                     & =(-1)^{f(0)}
    \begin{cases}
      \gate{H}_1\ket{+-} \\
      \gate{H}_1\ket{--}
    \end{cases}
    =(-1)^{f(0)}
    \begin{cases}
      \ket{0-} & \text{ si $f$ es constante}   \\
      \ket{1-} & \text{ si $f$ es equilibrada}
    \end{cases}
  \end{split}
\end{equation}

Por tanto, medir el primer cúbit en la base computacional da $0$ si $f$ es constante y $1$ si $f$ es equilibrada, con una única evaluación de $\gate{U}$.
